\chapter{Correctness, Validation, and Equivalence}
\label{chap:validation}

\section{Four Layers of Verification}
\label{sec:validation-layers}

% summary of the four layers; some chart

\section{Layer 1: Grounded Generation}
\label{sec:validation-grounding}

% pom.xml, .csproj, schema inspection, few-shot examples, how they are injected into the prompt; what the model sees and what it outputs; how the output is structured and validated against the schema; how the output is used in the next step

\section{Layer 2: Compilation and Execution in Sandboxes}
\label{sec:validation-sandbox}

% Daytona; how to build bash scripts run in Daytona sandboxes; transportation of files using base64; outputs of build tools, how long they are

\section{Layer 3: Semantic Equivalence}
\label{sec:validation-equivalence}


% explain entire DeepDiff algorithm using pseudocode; show JSON values; how they are compared; how the diff is tolerant to field order and rounding; what output to expect from DeepDiff; how it is passed to Judge
% reference ExeCoder a lot

\subsection{Tolerant comparison}

% how we defined tolerance threshold (LLM decides), but mostly JSONSerializer in services; why the JSONs have the same structure:
% differences in JSON serializers: Datetime, floating point, null, DateTime -> ZonedDateTime only, arrays/objects, field order

\subsection{The swapped-sorting robustness check}

% why we check firstSample-firstSample and lastSample-lastSample, and then firstSample-lastSample + lastSample-firstSample as well; threat to validity in terms of deterministic sorting

\section{Layer 4: The Independent Judge}
\label{sec:validation-judge}

% context it gets; output it generates; conditional edge after decides whether to go back to translation or END and show the user final translation

\section{Why the Result Can Be Trusted}
\label{sec:validation-trust}

% summary of the above; or omit
